\documentclass[11pt]{article}

\usepackage[margin=1in]{geometry}
\usepackage{amsmath,amssymb}
\usepackage{booktabs}
\usepackage{array}
\usepackage{xcolor}
\usepackage{hyperref}
\usepackage{tikz}
\usetikzlibrary{positioning}

\hypersetup{colorlinks=true,linkcolor=blue,citecolor=blue,urlcolor=blue}
\allowdisplaybreaks

\newcommand{\F}{\mathcal F}
\newcommand{\U}{\mathcal U}
\newcommand{\Pcal}{\mathcal P}
\newcommand{\SL}{\mathrm{SL}}
\newcommand{\HR}{\mathrm{HR}}
\newcommand{\zb}{\bar z}
\newcommand{\wb}{\bar w}
\newcommand{\hX}{\widehat X}
\newcommand{\eps}{\epsilon}

\title{NMRK and double-spacelike-collinear limits:\\
kinematics and hidden regions on equal footing}
\author{Draft note for internal discussion}
\date{21 July 2026}

\begin{document}
\maketitle

\section{Purpose and conventions}

This note compares two six-point asymptotic expansions using the same graphs,
the same external labels and the same final HRF certification criterion:
central next-to-multi-Regge kinematics (NMRK), with particular attention to
the surface \(w_M=z_M\), and the generic double-spacelike-collinear (DSC) limit of
Ref.~\cite{Duhr:2025lyg}.  Neither limit is defined by dual conformal
invariance.  Dual conformal invariance only reduces the variables on which a
particular observable can depend.

The presentation is deliberately not chronological.  We first define the two
limits symmetrically, then work through the common one-loop hexagon, and only
afterwards report and compare the other graphs.  Discovery-stage vectors are
kept distinct from certified region vectors.

All momenta are outgoing in algebraic formulae.  The physical sheet is a
\(2\to4\) process with particles 1 and 2 incoming.  Hence
\begin{equation}
 s_{12}>0,\qquad s_{34},s_{45},s_{56}>0,\qquad
 s_{16},s_{23}<0,
 \label{eq:physical-signs}
\end{equation}
where \(s_{16}\) and \(s_{23}\) become small in the DSC limit.  A region vector
always means
\begin{equation}
 x_i\longrightarrow\eta^{v_i}x_i,\qquad (v_0,\ldots,v_{N-1};1).
 \label{eq:vector-convention}
\end{equation}
Here \(\eta\) denotes the small parameter of the expansion under discussion:
the NMRK parameter in Sec.~2.2 or \(\varepsilon_D\) in Sec.~2.3.
The final component is displayed explicitly whenever there is a possible
normalization ambiguity.

\section{Kinematics: coordinates first, limits second}

Two notationally different objects were conflated in an earlier version of
this note: a coordinate system for generic six-point kinematics and a limit
taken inside that coordinate system.  They are separated here.  Subscripts
\(M\) and \(D\) distinguish the minimal-light-cone and Duhr--Venkata--Zhang
transverse variables.

\subsection{General minimal light-cone variables}

The minimal light-cone variables (MSLCV) follow
Refs.~\cite{Byrne:2022wdr,Byrne:2025phh}.  Set \(P=p_4^+\) and
\(Q=q_1^\perp\bar q_1^\perp\).  The longitudinal coordinates are defined by
\begin{equation}
p_3^+=PX_{34},\qquad p_4^+=P,\qquad
p_5^+=\frac{P}{X_{45}},\qquad
p_6^+=\frac{P}{X_{45}X_{56}},
\label{eq:nmrk-plus-definitions}
\end{equation}
and the transverse coordinates by
\begin{align}
p_4^\perp&=-\frac{q_1^\perp}{z_M-1},&
p_5^\perp&=\frac{z_Mq_1^\perp}{w_M(z_M-1)},\nonumber\\
p_6^\perp&=\frac{z_M(w_M-1)q_1^\perp}{w_M(z_M-1)},&
p_3^\perp&=-p_4^\perp-p_5^\perp-p_6^\perp,
\label{eq:nmrk-transverse-definitions}
\end{align}
with barred equations defining \(\bar z_M,\bar w_M\).  In complexified
kinematics the barred and unbarred coordinates may be manipulated as
independent algebraic variables.  On the real Lorentzian \(2\to4\) sheet
specified in Eq.~\eqref{eq:physical-signs}, however, the barred transverse
components are the complex conjugates of the unbarred ones.  Consequently,
\begin{equation}
 \bar z_M=z_M^*,\qquad \bar w_M=w_M^*.
 \label{eq:mslcv-reality}
\end{equation}
Thus the bars on the MSLCV variables used in the NMRK analysis below are
literal complex conjugation; algebraic independence is only a device for
complexification and analytic continuation.  Equations
\eqref{eq:nmrk-plus-definitions}--\eqref{eq:nmrk-transverse-definitions} are
a parametrization of generic kinematics; no Regge or collinear limit has yet
been taken.

The same coordinates can be reconstructed from invariants.  In the present
ordering,
\begin{equation}
 X_{34}=\frac{s_{13}}{s_{14}},\qquad
 X_{45}=\frac{s_{14}}{s_{15}},\qquad
 X_{56}=\frac{s_{15}}{s_{16}},\qquad
 Q=-\frac{s_{23}X_{34}X_{45}X_{56}}
 {1+X_{56}+X_{45}X_{56}+X_{34}X_{45}X_{56}}.
 \label{eq:mslcv-invariant-inversion}
\end{equation}
For completeness, put
\(u=z_M-1\), \(\bar u=\bar z_M-1\),
\(v=w_M-1\), and \(\bar v=\bar w_M-1\).  The transverse coordinates are
fixed, up to exchanging the two roots, by
\begin{align}
u\bar u&=\frac{s_{23}X_{34}}{s_{24}},&
u+\bar u&=\frac{X_{34}^2+u\bar u-s_{34}X_{34}u\bar u/Q}{X_{34}},
\nonumber\\
w_M\bar w_M&=\frac{X_{45}(1+u+\bar u+u\bar u)s_{24}}{s_{25}},&
v\bar v&=\frac{s_{26}}{X_{56}s_{25}},\qquad
v+\bar v=w_M\bar w_M-1-v\bar v.
\label{eq:mslcv-transverse-inversion}
\end{align}
Equations \eqref{eq:mslcv-invariant-inversion} and
\eqref{eq:mslcv-transverse-inversion} will provide the exact bridge from the
twistor parametrization below; they are valid before either limit is taken.

\subsection{Central NMRK and the additional \texorpdfstring{\(w_M=z_M\)}{wM=zM} surface}

Central NMRK is the scaling
\begin{equation}
 X_{34}=\frac{\hX_{34}}{\eta},\qquad
 X_{56}=\frac{\hX_{56}}{\eta},\qquad
 X_{45},Q,z_M,\bar z_M,w_M,\bar w_M=O(1),\qquad \eta\to0.
 \label{eq:nmrk-scaling}
\end{equation}
The positive HRF chart used below is
\begin{equation}
 Q,\hX_{34},X_{45},\hX_{56},K_z>0,\qquad
 K_z=(1-z_M)(1-\bar z_M)>0.
 \label{eq:nmrk-domain}
\end{equation}
The condition \(w_M=z_M\), \(\bar w_M=\bar z_M\), is a further restriction,
not part of the definition of NMRK.  For the planar ordering relevant to the
central-emission variables,
\begin{equation}
 1-u_A=\frac{X_{45}(w_M-z_M)(\bar w_M-\bar z_M)}
 {(1+X_{45})(w_M\bar w_M+X_{45}z_M\bar z_M)}.
 \label{eq:uA}
\end{equation}
Thus the extra surface is the distinguished cross-ratio boundary \(u_A=1\).

\subsection{DSC momentum twistors, labels and physical sheet}

We now define the variables actually used by the DSC runner.  Hats denote
the native cyclic labels of Ref.~\cite{Duhr:2025lyg}.  A convenient
projective gauge for the four undeformed twistors is
\begin{equation}
 Z_1=(1,0,0,0),\quad Z_4=(0,1,0,0),\quad
 Z_5=(1,1,1,0),\quad Z_6=(1,a,0,1),
 \label{eq:dsc-seed-twistors}
\end{equation}
where \(a\) is the remaining hard seed modulus.  With
\(\varepsilon_D\to0\), the double-collinear deformation is
\begin{align}
Z_2={}&Z_1+\varepsilon_D Z_4
-\varepsilon_D\tau_1\frac{\langle14\rangle}{\langle61\rangle}Z_6
-\varepsilon_D^2(1-z_D)\tau_2
 \frac{\langle14\rangle}{\langle45\rangle}Z_5,
\nonumber\\
Z_3={}&Z_4+\varepsilon_D Z_1
-\varepsilon_D\tau_2\frac{\langle14\rangle}{\langle45\rangle}Z_5
+\varepsilon_D^2\frac{(1-\bar z_D)\tau_1}{\bar z_D}
 \frac{\langle14\rangle}{\langle61\rangle}Z_6.
\label{eq:dsc-twistor-deformation}
\end{align}
The paper's momentum fractions obey
\(\tau_i=(1-\xi_i)/\xi_i\).  The variables \(z_D,\bar z_D\) in
Eq.~\eqref{eq:dsc-twistor-deformation} are not the MSLCV variables
\(z_M,\bar z_M\), despite the deliberately similar notation.

It is important to state the reality condition before continuing to the
DSC sheet.  In the standard real transverse-momentum realization of the
six-point \(2\to4\) multi-Regge region, the variables denoted by \(z_D\) and
\(\bar z_D\) form a complex-conjugate pair,
\begin{equation}
 \bar z_D=z_D^*.
 \label{eq:dsc-conjugate-branch}
\end{equation}
This is the branch on which the familiar MRK formulae are written in terms
of \(|z_D|\) and \(z_D/\bar z_D\).  Duhr--Venkata--Zhang reach the
double-spacelike-collinear region by analytic continuation from the timelike
Euclidean region into a specified Mandelstam region.  During this
continuation \(z_D\) and \(\bar z_D\) must be regarded as the two algebraic
roots, not as a conjugate pair at every intermediate point.

The change of reality is transparent from their cross-ratio variables.  If
\begin{equation}
 r_1=\frac{u_1}{1-u_2},\qquad r_3=\frac{u_3}{1-u_2},
\end{equation}
then their defining relations imply
\begin{equation}
 z_D\bar z_D=\frac{r_1}{r_3},\qquad
 z_D+\bar z_D=1+\frac{r_1-1}{r_3}.
 \label{eq:dsc-root-sum-product}
\end{equation}
Thus \(z_D\) and \(\bar z_D\) are the roots of a quadratic with real
coefficients whenever the continued ratios are real.  A negative quadratic
discriminant gives Eq.~\eqref{eq:dsc-conjugate-branch}; after the
discriminant changes sign the two roots are both real.  The latter is the
branch used by the DSC HRF chart below.

For clarity, the twistor-to-invariant operation used below is also fixed
here.  Write \(Z_i=(\lambda_i,\mu_i)\), take all labels cyclically, and set
\begin{equation}
 \widetilde\lambda_i=
 \frac{\langle i\,i{+}1\rangle\mu_{i-1}
 +\langle i{+}1\,i{-}1\rangle\mu_i
 +\langle i{-}1\,i\rangle\mu_{i+1}}
 {\langle i{-}1\,i\rangle\langle i\,i{+}1\rangle},
 \qquad
 s_{ij}^{(T)}=\langle ij\rangle[ji].
 \label{eq:twistor-to-invariants}
\end{equation}
Angle and square brackets are the two-by-two determinants of the
corresponding spinors.  Equations
\eqref{eq:dsc-seed-twistors}--\eqref{eq:twistor-to-invariants} therefore
determine every pair invariant rationally at finite \(\varepsilon_D\).

The native twistor labels are not the graph labels.  The order-aware map used
in the runner is
\begin{equation}
 p_1=\hat p_1,\quad p_2=\hat p_3,\quad p_3=\hat p_4,\quad
 p_4=\hat p_5,\quad p_5=\hat p_6,\quad p_6=\hat p_2.
 \label{eq:dsc-label-map}
\end{equation}
It maps the two small twistor pairs to the desired graph invariants
\(s_{23}\) and \(s_{16}\).  To keep the sign convention transparent, let
\(s_{ij}^{(T)}=\langle ij\rangle[ji]\) denote the invariants reconstructed
from the twistors in the seed convention above.  Their leading adjacent
values are
\begin{equation}
\begin{aligned}
s_{12}^{(T)}&=\frac{\tau_2}{a(1+\tau_1)(1+\tau_2)}+O(\varepsilon_D),&
s_{23}^{(T)}&=-\frac{\varepsilon_D^2\tau_1}{a\bar z_D}+O(\varepsilon_D^3),\\
s_{34}^{(T)}&=-\frac{1}{(a-1)(1+\tau_2)}+O(\varepsilon_D),&
s_{45}^{(T)}&=\frac1a,\\
s_{56}^{(T)}&=-\frac{\tau_1}{(a-1)(1+\tau_1)}+O(\varepsilon_D),&
s_{16}^{(T)}&=-\frac{\varepsilon_D^2\tau_2z_D}{a}+O(\varepsilon_D^3).
\end{aligned}
\label{eq:dsc-runner-leading-adjacent}
\end{equation}
Our graph-polynomial convention is
\begin{equation}
 s_{ij}=-s_{ij}^{(T)}.
 \label{eq:dsc-sign-convention}
\end{equation}
This is a single spinor-bracket/metric convention, not a separate sign choice
for each invariant.

A convenient \emph{real slice} of the physical \(2\to4\) sheet is
\begin{equation}
 a<0,\qquad -1<\tau_1<0,\qquad \tau_2<-1,\qquad
 z_D<0,\qquad \bar z_D<0.
 \label{eq:dsc-domain}
\end{equation}
On this slice \((s_{12}^{(T)},s_{23}^{(T)},s_{34}^{(T)},s_{45}^{(T)},
s_{56}^{(T)},s_{16}^{(T)})\) has signs
\((- ,+,-,-,-,+)\).  Equation \eqref{eq:dsc-sign-convention} therefore gives
exactly Eq.~\eqref{eq:physical-signs}.  The bar on \(\bar z_D\) is inherited
notation from the conjugate branch in Eq.~\eqref{eq:dsc-conjugate-branch}.
On the continued branch \eqref{eq:dsc-domain}, however, \(z_D\) and
\(\bar z_D\) are two real roots and need not be equal.  Requiring both roots
to be negative selects the sign-definite component used by HRF; it does not
reduce the kinematics to the diagonal \(z_D=\bar z_D\).  This is different
from Eq.~\eqref{eq:mslcv-reality}: the MSLCV variables remain complex
conjugates on the physical NMRK sheet considered in this note.

\subsection{Exact coordinate map and comparison of the two limiting paths}

There is no additional informal ``dictionary.''  At finite
\(\varepsilon_D\), reconstruct the fifteen \(s_{ij}\) using
Eqs.~\eqref{eq:dsc-seed-twistors}--\eqref{eq:dsc-sign-convention}, and insert
them in the invariant relations \eqref{eq:mslcv-invariant-inversion} and
\eqref{eq:mslcv-transverse-inversion}.  This is an exact, invariantly defined
map from \((a,\tau_1,\tau_2,z_D,\bar z_D,\varepsilon_D)\) to the MSLCV.

Expanding that exact map makes the limiting path explicit:
\begin{align}
X_{34}&=\frac{a-1}{1+\tau_2}+O(\varepsilon_D),&
X_{45}&=-\frac1a+O(\varepsilon_D),\nonumber\\
X_{56}&=\frac{a}{(a-1)(1+\tau_1)\tau_2z_D}\,
 \varepsilon_D^{-2}+O(\varepsilon_D^{-1}),&
Q&=\frac{\tau_1}{a\tau_2\bar z_D}\,
 \varepsilon_D^2+O(\varepsilon_D^3).
\label{eq:dsc-to-mslcv-longitudinal}
\end{align}
The approach of the transverse variables to the boundary is fixed by
\begin{align}
u\bar u&=\frac{(a-1)^2\tau_1}{a\tau_2\bar z_D}\,
 \varepsilon_D^2+O(\varepsilon_D^3),&
u+\bar u&=-\frac{(a-1)(\tau_1+a\tau_2\bar z_D)}
 {a\tau_2\bar z_D}\,\varepsilon_D+O(\varepsilon_D^2),\nonumber\\
v\bar v&=\frac{(a-1)^2\tau_1\tau_2z_D}{a}\,
 \varepsilon_D^2+O(\varepsilon_D^3),&
v+\bar v&=\frac{(a-1)(\tau_1+a\tau_2z_D)}{a}\,
 \varepsilon_D+O(\varepsilon_D^2).
\label{eq:dsc-to-mslcv-transverse}
\end{align}
Thus DSC has \(X_{34},X_{45}=O(1)\),
\(X_{56}=O(\varepsilon_D^{-2})\), \(Q=O(\varepsilon_D^2)\), and
\(z_M,\bar z_M,w_M,\bar w_M\to1\).  Central NMRK instead has both
\(X_{34}\) and \(X_{56}\) large while \(Q\) remains hard.  Generic DSC and
central NMRK with \(w_M=z_M\) are therefore distinct paths in the same
six-point kinematic space.  Their one-loop cancellation hypersurfaces are
nevertheless related, as shown explicitly in Sec.~\ref{sec:one-loop-common}.

\section{Common one-loop hexagon}
\label{sec:one-loop-common}

The graph and all label conventions are displayed in
Fig.~\ref{fig:one-loop-hexagon}.  The loop edges
\begin{equation}
 (5,6),(6,1),(1,2),(2,3),(3,4),(4,5)
 \quad\longleftrightarrow\quad
 x_0,x_1,x_2,x_3,x_4,x_5,
 \label{eq:hexagon-edge-map}
\end{equation}
respectively.  The external attachments are
\(p_1\to1\), \(p_2\to4\), \(p_3\to5\), \(p_4\to3\), \(p_5\to2\), and
\(p_6\to6\).

\begin{figure}[ht]
\centering
\begin{tikzpicture}[scale=1.15, every node/.style={font=\small}]
  \foreach \n/\ang in {1/150,2/90,3/30,4/-30,5/-90,6/-150}
    \coordinate (v\n) at (\ang:1.55);
  \draw (v1)--node[above left]{$x_2$}(v2)
        --node[above right]{$x_3$}(v3)
        --node[right]{$x_4$}(v4)
        --node[below right]{$x_5$}(v5)
        --node[below left]{$x_0$}(v6)
        --node[left]{$x_1$}(v1);
  \foreach \n in {1,...,6} \fill (v\n) circle (1.5pt);
  \node[above left=1pt of v1] {$1$};
  \node[above=2pt of v2] {$2$};
  \node[above right=1pt of v3] {$3$};
  \node[below right=1pt of v4] {$4$};
  \node[below=2pt of v5] {$5$};
  \node[below left=1pt of v6] {$6$};
  \draw (v1)--++(150:0.75) node[above left]{$p_1$};
  \draw (v2)--++(90:0.75) node[above]{$p_5$};
  \draw (v3)--++(30:0.75) node[above right]{$p_4$};
  \draw (v4)--++(-30:0.75) node[below right]{$p_2$};
  \draw (v5)--++(-90:0.75) node[below]{$p_3$};
  \draw (v6)--++(-150:0.75) node[below left]{$p_6$};
\end{tikzpicture}
\caption{The one-loop hexagon, with the graph vertices, external momenta and
Schwinger parameters used by HRF.}
\label{fig:one-loop-hexagon}
\end{figure}

Before either kinematic expansion, the Symanzik polynomials are
\begin{equation}
 \U=x_0+x_1+x_2+x_3+x_4+x_5.
\end{equation}
and
\begin{align}
\F={}&s_{16}x_0x_2
+(-s_{14}+s_{23}-s_{45}-s_{46})x_0x_3+s_{23}x_0x_4
\nonumber\\
&+(-s_{14}-s_{16}+s_{23}-s_{45}-s_{46}-s_{56})x_1x_3
\nonumber\\
&+(-s_{16}+s_{23}-s_{46}-s_{56})x_1x_4
\nonumber\\
&+(s_{12}-s_{34}-s_{35}-s_{45}-s_{46}-s_{56})x_1x_5
\nonumber\\
&+s_{45}x_2x_4+(s_{16}-s_{23}-s_{34}-s_{35})x_2x_5
\nonumber\\
&+(-s_{14}-s_{34}-s_{45}-s_{46})x_3x_5.
\label{eq:hexagon-F}
\end{align}
The same \(\U\), \(\F\), graph labels and HRF implementation are used for
both expansions.  Only the substitution for the invariants \(s_{ij}\)
changes.

\subsection{Certification rule after asymptotic-order alignment}

This subsection fixes the notation used in the two examples.  Write the full
kinematically expanded polynomial as
\begin{equation}
 \F(\eta,\mathbf x)=\sum_m \eta^{a_m}c_m\mathbf x^{\mathbf r_m}.
 \label{eq:full-laurent-F}
\end{equation}
The algorithm has three logically separate stages.

\paragraph{1. Asymptotic-order alignment (face discovery).}
A trial face vector \(f\) assigns the total weight
\(a_m+\mathbf r_m\mathbin{\cdot}f\) to each monomial.  The monomials of
minimum weight define an order-aligned face \(\F^{(f)}\).  This operation
exposes monomials that occur at different fixed-\(\mathbf x\) orders in
\(\eta\), but acquire a common effective order after the parameter scaling
\(x_e\mapsto\eta^{f_e}x_e\), so that they can participate in the same
cancellation sector.  Because \(\F\) is
homogeneous in the Schwinger parameters, replacing
\(f\to f+c(1,\ldots,1)\) does not change which monomials lie on the face.
Thus \(f\) is only a discovery-stage representative, not yet a physical
region vector.

\paragraph{2. Ordinary HRF on the selected face.}
Once \(\F^{(f)}\) has been selected, the usual HRF analysis is performed just
as it would be without asymptotic-order alignment.  HRF finds non-monomial cancellation
factors \(L_i(\mathbf x)\), forms their ideal
\(I=\langle L_1,\ldots,L_r\rangle\), and decomposes
\begin{equation}
 \F^{(f)}=\F_{\SL}+\F_{\rm obs},\qquad \F_{\SL}\in I.
 \label{eq:hrf-decomposition}
\end{equation}
It then finds a cancellation scaling \(h\) for which all monomials of
\(\F_{\SL}\) have the common weight \(W_{\SL}\), while the first surviving
terms of \(\U+\F_{\rm obs}\) have weight \(W_{\HR}>W_{\SL}\).  We call \(h\)
the \emph{HRF face scaling}.  The displayed sum \(f+h\) is useful during
discovery, but it inherits the uniform ambiguity of \(f\) and is not
automatically the final vector.

\paragraph{3. Certification in the original coordinates.}
The algorithm returns to the full \(\Pcal(\eta,\mathbf x)=\U+\F\) and retains
every occupied weight layer.  Each \(\F\)-layer is resolved in the filtration
\(I^q/I^{q+1}\).  A layer in \(I^q\) cancels to order \(q\) at the common
pinch; an absent power of \(\eta\) is not counted as a cancellation layer.
The first resolved nonvanishing layer fixes \(W_{\HR}\).  The simultaneous
balance with \(\U\) fixes the previously ambiguous uniform component and,
with the convention in Eq.~\eqref{eq:vector-convention}, produces the final
original-coordinate vector \(v\).

If the resolved \(W_{\HR}\) monomials already define a unique lower facet of
the full Lee--Pomeransky polynomial, this gives a direct certificate.  If the
ideal-layer system is nonlinear, singular or non-unique, a local dissection
provides the necessary and sufficient facet test and its pullback fixes
\(v\).  No space-time dimensional information enters any of these decisions.

\subsection{NMRK with \texorpdfstring{\(w_M=z_M\)}{wM=zM}}

The cancellation surface factorizes as
\begin{align}
 L_1^{\rm NMRK}&=(1+X_{45})x_4-\hX_{34}X_{45}x_5,\nonumber\\
 L_2^{\rm NMRK}&=(1+X_{45})x_2-K_zX_{45}\hX_{56}x_1.
 \label{eq:nmrk-factors}
\end{align}
Up to a nonzero kinematic factor,
\begin{equation}
 \F_{\SL}^{\rm NMRK}=Q\hX_{34}\hX_{56}
 L_1^{\rm NMRK}L_2^{\rm NMRK}.
 \label{eq:nmrk-fsl}
\end{equation}
For generic NMRK the corresponding bilinear coefficient determinant is
proportional to
\((w_M-z_M)(\bar w_M-\bar z_M)\); it is a near miss.  The surface
\(w_M=z_M\) makes the
factorization exact.

For this case the discovery and certification data are
\begin{align}
 f&=(-1,-1,-2,-1,-2,-1),\nonumber\\
 h&=(-1,-1,-1,-1,-1,-1),\nonumber\\
 v_{\rm NMRK}&=f+h=(-2,-2,-3,-2,-3,-2),
 \qquad (v_{\rm NMRK};1),
 \label{eq:nmrk-vector}
\end{align}
with
\begin{equation}
 W_{\SL}=-4,\qquad W_{\HR}=-3,\qquad
 W_{\HR}-W_{\SL}=1.
 \label{eq:nmrk-gap}
\end{equation}
Here the direct composition is already the unique lower-facet normal.  At
\(W_{\HR}\), the active \(\U\) contribution is \(x_2+x_4\), and the active
\(\F\) contribution is
\begin{equation}
 -Q\hX_{34}X_{45}\hX_{56}
 (K_zx_0x_2+K_zx_0x_4+\hX_{34}x_3x_5
 +K_z\hX_{56}x_1x_3).
 \label{eq:nmrk-whr}
\end{equation}
These monomials, together with the resolved promoted face, certify a scaleful
region.

\subsection{Generic DSC}

The DSC cancellation factors are
\begin{equation}
 L_1^{\rm DSC}=(1+\tau_2)x_4+x_5,\qquad
 L_2^{\rm DSC}=\tau_1x_1+(1+\tau_1)x_2,
 \label{eq:dsc-factors}
\end{equation}
and
\begin{equation}
 \F_{\SL}^{\rm DSC}=
 -(a-1)a^2\bar z_D^{,2}L_1^{\rm DSC}L_2^{\rm DSC}.
 \label{eq:dsc-fsl}
\end{equation}
These factors make the contrast with the crossed double-timelike-collinear
(DTC) domain immediate.  In the DSC domain
\(-1<\tau_1<0\) and \(\tau_2<-1\), each factor has terms of opposite sign
and admits a zero with positive Feynman parameters:
\begin{equation}
 L_1^{\rm DSC}=0:\quad
 x_5=-(1+\tau_2)x_4>0,
 \qquad
 L_2^{\rm DSC}=0:\quad
 x_2=-\frac{\tau_1}{1+\tau_1}x_1>0.
 \label{eq:dsc-positive-pinch}
\end{equation}
The two equations can therefore hold simultaneously in the interior of the
integration domain.  For two final-state timelike splittings, by contrast,
the momentum fractions satisfy \(0<\xi_i<1\), and hence
\begin{equation}
 \tau_i=\frac{1-\xi_i}{\xi_i}>0.
 \label{eq:dtc-tau-domain}
\end{equation}
It follows that both \(L_1^{\rm DSC}\) and \(L_2^{\rm DSC}\) are strictly
positive for positive \(x_i\).  Crossing thus preserves the algebraic form
of this candidate generator but removes its common positive Landau pinch.
This excludes the timelike continuation of the DSC hidden region; by itself
it does not exclude a different cancellation generator in the DTC limit.

The discovery-stage vectors are
\begin{align}
 f&=(-2,-2,-2,-1,-2,-2),\nonumber\\
 h&=(-1,-1,-1,-1,-1,-1),\nonumber\\
 f+h&=(-3,-3,-3,-2,-3,-3).
 \label{eq:dsc-naive}
\end{align}
The last line is not the final region vector.  The full nonzero coefficient
layers satisfy
\begin{equation}
 \operatorname{ord}_I\F_{-4}=2,\qquad
 \operatorname{ord}_I\F_{-3}=1,\qquad
 \operatorname{ord}_I\F_{-2}=0,\qquad
 I=\langle L_1^{\rm DSC},L_2^{\rm DSC}\rangle.
\label{eq:dsc-ideal-orders}
\end{equation}
Here \(\F_k\) is the sum of monomials of total weight \(k\), and
\(\operatorname{ord}_I\F_k=q\) means
\(\F_k\in I^q\) but \(\F_k\notin I^{q+1}\).
Thus the \(-3\) layer is present but cancels to first order on the common
pinch.  Solving the ideal-jet balance gives transverse weights
\(q_1=q_2=1\), and \(\U\) fixes the uniform component.  The certified answer is
\begin{equation}
 v_{\rm DSC}=(-2,-2,-2,0,-2,-2),\qquad (v_{\rm DSC};1),
 \label{eq:dsc-vector}
\end{equation}
with
\begin{equation}
 W_{\SL}=-4,\qquad W_{\HR}=-2,\qquad
 W_{\HR}-W_{\SL}=2.
 \label{eq:dsc-gap}
\end{equation}
This is an original-coordinate ideal-jet certificate.  Dissection pulls back
to the same vector and explains the previously observed factor of two.

\subsection{What is common and what is different}

After the separate kinematic rescalings have been made, the coefficients in
the two factorized \(\F_{\SL}\) polynomials are identified by
\begin{equation}
 \frac{K_zX_{45}\hX_{56}}{1+X_{45}}=-\frac{\tau_1}{1+\tau_1},
 \qquad
 \frac{\hX_{34}X_{45}}{1+X_{45}}=-\frac{1}{1+\tau_2}.
 \label{eq:factor-map}
\end{equation}
This is a map of the cancellation factors on the two selected faces, not a
claim that the full kinematic limits coincide.  Consequently the factorized
\(W_{\SL}\) support is the same.  The expansions are nevertheless distinct because the
next occupied layers differ.  NMRK with \(w_M=z_M\) has a depth-one hierarchy;
generic DSC has an additional first-order cancellation layer and hence a
depth-two hierarchy.  This is why the two final vectors are related but not
equal.

\begin{center}
\begin{tabular}{@{}lccc@{}}
\toprule
Expansion & certified \(v\) & \((W_{\SL},W_{\HR})\) & gap\\
\midrule
NMRK with \(w_M=z_M\) & \((-2,-2,-3,-2,-3,-2)\) & \((-4,-3)\) & 1\\
generic DSC & \((-2,-2,-2,0,-2,-2)\) & \((-4,-2)\) & 2\\
\bottomrule
\end{tabular}
\end{center}

\section{Results within each limit}

\subsection{NMRK results}

The one-loop result above and all three two-loop controls have now been
passed through the current final audit.  ``Full'' means that no Feynman
parameter is fixed to zero; ``boundary'' denotes a certified scaling vector
on a one-edge boundary stratum.  At two loops, a certified vector is not yet
identified with a distinct physical hidden region: vectors must also be
compared at the level of their polynomial generator and local facet.
Every vector in the tables below has last component \(1\).  When a vector is
called ``relative,'' its uniform representative is chosen so that
\(\max_i v_i=0\); this does not suppress or alter the final component.

\begin{center}
\begin{tabular}{@{}lccc@{}}
\toprule
Graph & generic NMRK & certified vectors at \(w_M=z_M\) & support split\\
\midrule
one-loop hexagon & no HR & 1 certified & 1 full\\
planar hexbox & no staged candidate & 11 certified & 7 full + 4 boundary\\
non-planar hexbox & no staged candidate & 3 certified & 1 full + 2 boundary\\
hexagon-pentagon & no staged candidate & no staged candidate & --\\
\bottomrule
\end{tabular}
\end{center}

\subsubsection*{Planar-hexbox generator multiplicity}

A cancellation hypersurface is defined by non-monomial polynomial factors.
No individual Feynman parameter \(x_i\) is a cancellation hypersurface, and
no monomial factor is counted as an HR generator.  For the seven
full-support planar-hexbox vectors, introduce the non-monomial linear
polynomials
\begin{align}
A&=(1+X_{45})x_3-K_zX_{45}\hX_{56}x_1,&
B&=(1+X_{45})x_3-K_zX_{45}\hX_{56}x_2,\nonumber\\
C&=(1+X_{45})x_3-K_zX_{45}\hX_{56}(x_1+x_2),&
D&=(1+X_{45})x_5-\hX_{34}X_{45}x_6,\nonumber\\
E&=(1+X_{45})x_5-\hX_{34}X_{45}x_7,&
F&=(1+X_{45})x_5-\hX_{34}X_{45}(x_6+x_7).
\label{eq:planar-generator-factors}
\end{align}
Every factor in Eq.~\eqref{eq:planar-generator-factors} contains at least two
Schwinger monomials.  Up to factor ordering and overall signs, the seven
vectors form five generator classes,
\begin{equation}
G_1=AD,\qquad G_2=CD,\qquad G_3=AF,\qquad
G_4=AE,\qquad G_5=BD.
\label{eq:planar-generator-classes}
\end{equation}
The detailed assignment is
\begin{center}
\small
\begin{tabular}{@{}ccl@{}}
\toprule
row & generator & relative certified vector\\
\midrule
1 & \(G_1\) & \((-2,-1,-1,-2,-1,-2,-1,0,-1)\)\\
2 & \(G_2\) & \((-1,-1,-1,-2,-1,-2,-1,0,-2)\)\\
3 & \(G_3\) & \((-1,-1,0,-2,-1,-2,-1,-1,-2)\)\\
4 & \(G_1\) & \((-1,-1,0,-2,-1,-2,-1,0,-2)\)\\
5 & \(G_1\) & \((-1,-1,0,-2,-1,-2,-1,0,-1)\)\\
6 & \(G_4\) & \((-1,-1,0,-2,-1,-2,0,-1,-2)\)\\
7 & \(G_5\) & \((-1,0,-1,-2,-1,-2,-1,0,-2)\)\\
\bottomrule
\end{tabular}
\end{center}
Rows 1, 4 and 5 therefore have exactly the same polynomial generator
\(G_1=AD\), but different certified vectors.  The present lower-facet audit
proves that each vector is scaleful; it does not prove that these three
vectors define three physically distinct hidden regions.  That distinction
requires a common local-coordinate analysis around \(A=D=0\), followed by a
single dissection/facet enumeration and a comparison of the active monomial
sets and normal cones.  The count ``seven'' must consequently be read as
seven certified full-support vectors, not seven established physical HRs.

\subsection{DSC results}

The saved DSC scans have all been passed through the current structural
audit.  ``Staged'' counts face presentations before the final
scalefulness/layer test; ``certified'' counts accepted structural vectors.

\begin{center}
\begin{tabular}{@{}lrrl@{}}
\toprule
Graph & staged presentations & certified & certified vector(s)\\
\midrule
one-loop hexagon & 1 structural class & 1 & \((-2,-2,-2,0,-2,-2)\)\\
planar hexbox & 156 & 1 & \((0,-2,0,-2,0,-2,-2,0,0)\)\\
non-planar hexbox & 61 & 0 & --\\
hexagon-pentagon & 0 & 0 & --\\
\bottomrule
\end{tabular}
\end{center}
For the planar hexbox, 27 structural representatives were audited: 26 fail
the resolved \(W_{\HR}\)-facet test and one has an accepted layered common
pinch facet.  The non-planar hexbox has six structural representatives, all
rejected by the resolved facet test.

\section{Graph-by-graph comparison}

The one-loop hexagon shows the closest relationship.  Both limits reach the
same factorized singular hypersurface, but DSC contains one extra occupied
cancellation layer.  This changes the hierarchy gap from 1 to 2 and changes
the final original-coordinate vector.

For the planar hexbox, NMRK with \(w_M=z_M\) has seven full-support certified
vectors and four certified vectors on one-edge boundary strata.  All have
hierarchy gap 1 and direct resolved-facet certificates.  The seven interior
vectors represent only five polynomial-generator classes, and three of them
share \(G_1=AD\), as detailed above.  Generic DSC instead has one
full-support certified vector,
\begin{equation}
v_{\rm DSC}^{\rm P}=(0,-2,0,-2,0,-2,-2,0,0),\qquad
(W_{\SL},W_{\HR})=(-4,-2),
\end{equation}
with gap 2 and a layered common-pinch dissection certificate.  None of the
seven full-support NMRK relative vectors equals this DSC vector.  This proves
that the displayed scaling vectors differ, but it does not by itself count
distinct physical regions.  A generator-level translation and common local
dissection are required before deciding how many planar-hexbox HRs the two
limits share.

For the non-planar hexbox, NMRK with \(w_M=z_M\) has one full-support certified
vector and two certified vectors on one-edge boundary strata, again with gap
1.  The full-support relative vector is
\begin{equation}
(-1,-1,-2,-1,-2,-1,-1,0,-2),
\end{equation}
whereas generic DSC has no certified region.  The hexagon-pentagon is a
common negative control: neither expansion has even a staged candidate.

\section{Reproducibility}

The calculations are separated by expansion:
\begin{itemize}
 \item \path{06_NMRK_wz_AsymptoticOrderAlignment_Checks.nb} contains only NMRK inputs and
 certificates;
 \item \path{07_DSC_AsymptoticOrderAlignment_Checks.nb} contains only DSC inputs and
 certificates;
 \item \path{HRF_RunNMRKMultiGraphAsymptoticOrderAlignment.wl} reruns the NMRK two-loop
 graphs with the current final audit;
 \item \path{HRF_RunGenericDSCAsymptoticOrderAlignment.wl} reruns the DSC graphs.
\end{itemize}
This document is the only layer that translates and compares the two
expansions.

\begingroup
\small
\begin{thebibliography}{9}
\bibitem{Byrne:2022wdr}
E.~P.~Byrne, V.~Del Duca, L.~J.~Dixon, E.~Gardi and J.~M.~Smillie,
``One-loop central-emission vertex for two gluons in
\(\mathcal N=4\) super Yang--Mills theory,'' JHEP \textbf{08} (2022) 271,
\href{https://arxiv.org/abs/2204.12459}{arXiv:2204.12459}.

\bibitem{Byrne:2025phh}
E.~P.~Byrne, V.~Del Duca, E.~Gardi, Y.~Mo and J.~M.~Smillie,
``Regge factorization of tree-level QCD amplitudes using a minimal set of
lightcone variables,'' JHEP \textbf{03} (2026) 170,
\href{https://arxiv.org/abs/2506.10644}{arXiv:2506.10644}.

\bibitem{Gardi:2024axt}
E.~Gardi, F.~Herzog, S.~Jones and Y.~Ma,
``Dissecting polytopes: Landau singularities and asymptotic expansions in
\(2\to2\) scattering,'' JHEP \textbf{08} (2024) 127,
\href{https://arxiv.org/abs/2407.13738}{arXiv:2407.13738}.

\bibitem{Duhr:2025lyg}
C.~Duhr, A.~Venkata and C.~Zhang,
``Double Spacelike Collinear Limits from Multi-Regge Kinematics,''
Phys. Rev. Lett. \textbf{135} (2025) 241601,
\href{https://arxiv.org/abs/2507.05355}{arXiv:2507.05355}.
\end{thebibliography}
\endgroup

\end{document}
